All figures and derived data products in this work are generated by a
single reproducible pipeline driven by a JSON configuration file. The
pipeline executes eleven sequential stages, each defined by a working
directory, a set of commands, and a list of expected output files. The
pipeline halts on first failure, and output files from earlier stages
are available as resolved variables in later stages. The stages
proceed as follows.

\begin{enumerate}

\item Fit the six-parameter flare frequency distribution (FFD) model
of \citet{Davenport19} to the \emph{Kepler} ensemble data using
\texttt{emcee} \citep{ForemanMackey13}, producing posterior samples
and a corner plot.

\item Compare the FFD slope and intercept across stellar age using
the fitted model.

\item Identify flares in the \emph{TESS} two-minute cadence
lightcurves of GJ~1132 and plot the three detected events.

\item Construct a comprehensive FFD comparison among GJ~1132, field
M~dwarfs from \citet{Hawley14}, and open cluster stars from
\citet{Ilin19}.

\item Derive the present-day XUV luminosity distribution from the
\citet{Ribas05} scaling relation using Monte Carlo sampling of the
measured bolometric luminosity and activity ratio.

\item Compute the stellar age distribution from the
\citet{Engle24} gyrochronology--activity relation via Monte Carlo
sampling.

\item Determine the maximum likelihood parameter set with
\texttt{MaxLEV} using differential evolution, then train a Gaussian
process surrogate for the log-likelihood with \texttt{alabi}
\citep{Birky25}. Kernel and scaler hyperparameters are selected by
grid search (Table~\ref{tab:grid_search}), followed by 500 iterations
of active learning. Posterior samples are drawn with four independent
samplers---\texttt{emcee}, \texttt{dynesty} \citep{Speagle20},
\texttt{PyMultiNest} \citep{Buchner14}, and \texttt{UltraNest}
\citep{Buchner21}---and compared in a single corner plot.

\item Run convergence testing with \texttt{vconverge} for two XUV
evolution models: \citet{Engle24} with the \citet{Barnes20} power-law
decay, and \citet{Ribas05} with the same decay law. Each model
iterates batches of \texttt{VPLanet} \citep{Barnes20} simulations
until the output distributions satisfy a two-sample
Kolmogorov--Smirnov convergence criterion, then plots the resulting
XUV luminosity evolution envelopes.

\item Compute the cumulative XUV fluence distributions from the
converged simulation ensembles and locate GJ~1132\,b on the cosmic
shoreline diagram.

\item Quantify the relative contribution of each uncertain input
parameter to the total variance in cumulative XUV fluence.

\item Compare the \emph{Kepler}-predicted and \emph{TESS}-measured
FFD parameters ($\alpha$, $\beta$) at the inferred age of GJ~1132.

\end{enumerate}

The pipeline is implemented in \texttt{Python} and orchestrated by a
generic driver script (\texttt{director.py}) that reads the stage
definitions from a JSON file, executes each stage in a subprocess with
prefixed logging, and verifies that all expected output files are
produced. The full pipeline, including all simulations, completes in
approximately six hours on eight cores. All source code and
configuration files needed to reproduce these results are available at
\dataset[doi]{URL}.
